% !TeX root = booklet-03.tex
مقدمه بخش \secnum{۲.۵} در صفحه چاپی ۸۷، بخش \secnum{۲.۵.۱}
در صفحات ۸۸ و ۸۹، جدول‌های \secnum{۲.۲} و \secnum{۲.۳} و شکل‌های
\secnum{۲.۲۵} تا \secnum{۲.۲۷} منابع بخش اول‌اند.
بخش \secnum{۲.۵.۲} در صفحه ۹۰ و ابتدای ۹۱، روابط
\secnum{۲.۵۴} تا \secnum{۲.۵۷} و شکل \secnum{۲.۲۸} مبنای بحث
نسبت قدرت هستند. مدل جریان \lr{Long-channel} از بخش
\secnum{۲.۲}، حد سرعت از \secnum{۲.۴.۱} و اثر
\lr{CLM} از بخش \secnum{۲.۴.۲} به کار رفته‌اند.
ساختار گیت‌های \lr{NAND} و \lr{NOR} با بخش \secnum{۱.۴}
در صفحات ۹ و ۱۱ تطبیق داده شده است.
بخش \secnum{۲.۵.۳} در صفحات ۹۱ و ۹۲، شکل‌های
\secnum{۲.۲۹} و \secnum{۲.۳۰} و روابط \secnum{۲.۵۸} و
\secnum{۲.۵۹} مبنای بحث نویز هستند.
بخش \secnum{۲.۵.۴} در صفحات ۹۲ و آغاز ۹۳ و شکل
\secnum{۲.۳۱} مبنای تحلیل انتقال ولتاژ است. رابطه اثر بدنه
از بخش \secnum{۲.۴.۳.۱} در صفحه ۷۹ به کار رفته است.
صورت اصلی تمرین‌های \secnum{۲.۱۶}، \secnum{۲.۱۷} و \secnum{۲.۲۰} و اتصال‌های شکل
\secnum{۲.۳۵} در صفحه ۹۷ بررسی شده‌اند.
شماره صفحه فایل \lr{PDF} برابر شماره چاپی به‌اضافه ۲۷ است.
